\documentclass[11pt,a4paper]{article}

% ============================================================
% SCT Theory — Prediction 02: Running Constants with Fixed ICs
% Status: UNTESTED
% ============================================================

\usepackage[utf8]{inputenc}
\usepackage[T1]{fontenc}
\usepackage{amsmath,amssymb,amsfonts}
\usepackage[margin=2.5cm]{geometry}
\usepackage{hyperref}
\usepackage{booktabs}
\usepackage{xcolor}
\usepackage{fancyhdr}
\usepackage{enumitem}

\definecolor{untested}{RGB}{180,120,0}
\definecolor{sctblue}{RGB}{0,51,153}

\pagestyle{fancy}
\fancyhf{}
\fancyhead[L]{\small\textsc{SCT Theory --- Prediction Card}}
\fancyhead[R]{\small\textsc{P-02}}
\fancyfoot[C]{\thepage}

\title{%
  \textbf{SCT Prediction 02:}\\[6pt]
  \Large Running of Curvature-Squared Couplings\\
  with Fixed Initial Conditions\\[4pt]
  \large $c_1^{\text{SCT}}(\mu) = \dfrac{f_4}{480\pi^2} - 2b_1\ln\!\left(\dfrac{\Lambda}{\mu}\right)$
}
\author{David Alfyorov}
\date{March 2026}

\begin{document}
\maketitle


\begin{center}
\small\textit{Credits: Aliaksandr Samatyia contributed research-assistance and workflow support.}
\end{center}

\begin{center}
  \colorbox{untested!15}{\parbox{0.6\textwidth}{\centering
    \textbf{\color{untested} STATUS: UNTESTED}\\[2pt]
    \small Current experiments cannot probe curvature-squared running
  }}
\end{center}

\vspace{1em}

% ============================================================
\section{Source}
% ============================================================

\begin{itemize}[leftmargin=2em]
  \item \textbf{Derived from:} Combination of:
    \begin{enumerate}
      \item SCT Axiom~1 (Spectral Action Principle), which fixes the bare values of $c_1$ and $c_2$ at the cutoff scale $\Lambda$ (see Prediction P-01).
      \item Standard renormalization group equations for higher-derivative gravity (Julve--Tonin, 1978; Fradkin--Tseytlin, 1982; Avramidi--Barvinsky, 1985).
    \end{enumerate}
  \item \textbf{Derivation chain:}
    \begin{enumerate}
      \item The spectral action fixes boundary conditions at $\mu = \Lambda$:
        \begin{align}
          c_1(\Lambda) &= \frac{f_4}{480\pi^2} \label{eq:c1-boundary}\\
          c_2(\Lambda) &= -\frac{f_4}{160\pi^2} \label{eq:c2-boundary}
        \end{align}
      \item The one-loop renormalization group equations for the curvature-squared couplings in the gravitational effective action are:
        \begin{align}
          \mu\frac{d c_1}{d\mu} &= -2b_1 \label{eq:beta-c1}\\
          \mu\frac{d c_2}{d\mu} &= -2b_2 \label{eq:beta-c2}
        \end{align}
        where $b_1, b_2$ are the one-loop beta function coefficients.
      \item Integrating~\eqref{eq:beta-c1}--\eqref{eq:beta-c2} from $\Lambda$ down to $\mu$:
        \begin{align}
          c_1^{\text{SCT}}(\mu) &= \frac{f_4}{480\pi^2} - 2b_1\ln\!\left(\frac{\Lambda}{\mu}\right) \label{eq:c1-running}\\
          c_2^{\text{SCT}}(\mu) &= -\frac{f_4}{160\pi^2} - 2b_2\ln\!\left(\frac{\Lambda}{\mu}\right) \label{eq:c2-running}
        \end{align}
    \end{enumerate}
  \item \textbf{Key references:}
    \begin{itemize}
      \item J.~Julve, M.~Tonin, \textit{Nuovo Cim.\ B}~\textbf{46}, 137 (1978).
      \item E.~S.~Fradkin, A.~A.~Tseytlin, \textit{Nucl.\ Phys.\ B}~\textbf{201}, 469 (1982).
      \item I.~G.~Avramidi, A.~O.~Barvinsky, \textit{Phys.\ Lett.\ B}~\textbf{159}, 269 (1985).
      \item A.~O.~Barvinsky, G.~A.~Vilkovisky, \textit{Phys.\ Rep.}~\textbf{119}, 1 (1985).
    \end{itemize}
\end{itemize}

% ============================================================
\section{Observable Quantity}
% ============================================================

\subsection{The Prediction}

The running curvature-squared couplings with \textbf{fixed initial conditions} determined by the spectral geometry:

\begin{equation}
  \boxed{
    c_1^{\text{SCT}}(\mu) = \frac{f_4}{480\pi^2} - 2b_1\ln\!\left(\frac{\Lambda}{\mu}\right)
  }
  \label{eq:prediction}
\end{equation}

\begin{equation}
  \boxed{
    c_2^{\text{SCT}}(\mu) = -\frac{f_4}{160\pi^2} - 2b_2\ln\!\left(\frac{\Lambda}{\mu}\right)
  }
  \label{eq:prediction2}
\end{equation}

\subsection{What Makes This a Prediction}

In standard EFT of gravity, the renormalization group equations~\eqref{eq:beta-c1}--\eqref{eq:beta-c2} are the same, but the initial conditions $c_1(\mu_0)$ and $c_2(\mu_0)$ at some reference scale $\mu_0$ are \textbf{free parameters} that must be measured experimentally. The full running is:
\begin{equation}
  c_i^{\text{EFT}}(\mu) = c_i(\mu_0) - 2b_i\ln\!\left(\frac{\mu_0}{\mu}\right)
  \label{eq:eft-running}
\end{equation}
with two undetermined constants $c_1(\mu_0), c_2(\mu_0)$.

In SCT, these constants are \textbf{replaced by computable quantities}:
\begin{equation}
  c_1(\mu_0) \to \frac{f_4}{480\pi^2} - 2b_1\ln\!\left(\frac{\Lambda}{\mu_0}\right), \qquad
  c_2(\mu_0) \to -\frac{f_4}{160\pi^2} - 2b_2\ln\!\left(\frac{\Lambda}{\mu_0}\right)
\end{equation}
The only remaining free parameter is $f_4$, which is a moment of the cutoff function:
\begin{equation}
  f_4 = \int_0^\infty f(u)\,u\,du
  \label{eq:f4-definition}
\end{equation}
This is fixed once the cutoff function $f$ and the spectral triple $(A, H, D)$ are specified.

\subsection{Beta Function Coefficients}

The one-loop beta function coefficients for pure gravity (no matter fields):
\begin{align}
  b_1^{\text{pure grav}} &= \frac{1}{(4\pi)^2}\cdot\frac{133}{10} \\
  b_2^{\text{pure grav}} &= -\frac{1}{(4\pi)^2}\cdot\frac{196}{45}
\end{align}

With $N_s$ real scalars, $N_f$ Dirac fermions, and $N_v$ vectors:
\begin{align}
  b_1 &= \frac{1}{(4\pi)^2}\left(\frac{133}{10} + N_s\cdot\frac{1}{2} + N_f\cdot(-1) + N_v\cdot\frac{11}{2}\right) \\
  b_2 &= \frac{1}{(4\pi)^2}\left(-\frac{196}{45} + N_s\cdot\frac{1}{60} + N_f\cdot\frac{1}{10} + N_v\cdot\frac{31}{180}\right)
\end{align}

For the Standard Model matter content ($N_s = 4, N_f = 45, N_v = 12$):
\begin{align}
  b_1^{\text{SM}} &= \frac{1}{(4\pi)^2}\left(\frac{133}{10} + 2 - 45 + 66\right) = \frac{1}{(4\pi)^2}\cdot\frac{363}{10} \\
  b_2^{\text{SM}} &= \frac{1}{(4\pi)^2}\left(-\frac{196}{45} + \frac{4}{60} + \frac{45}{10} + \frac{12\cdot 31}{180}\right) = \frac{1}{(4\pi)^2}\cdot\frac{53}{9}
\end{align}

% ============================================================
\section{SCT vs.\ Standard EFT}
% ============================================================

\begin{center}
\renewcommand{\arraystretch}{1.4}
\begin{tabular}{p{4.5cm}cc}
  \toprule
  & \textbf{SCT} & \textbf{Standard EFT} \\
  \midrule
  RG equations & Same & Same \\
  Beta functions $b_1, b_2$ & Same & Same \\
  Initial condition $c_1(\Lambda)$ & $f_4/(480\pi^2)$ & Free parameter \\
  Initial condition $c_2(\Lambda)$ & $-f_4/(160\pi^2)$ & Free parameter \\
  Ratio at $\mu = \Lambda$ & $-1/3$ (fixed) & Arbitrary \\
  Free parameters & 1 ($f_4$) & 2 ($c_1(\mu_0), c_2(\mu_0)$) \\
  Predictive power & Higher & Lower \\
  \bottomrule
\end{tabular}
\end{center}

\vspace{0.5em}
\noindent
\textbf{Key distinction:} SCT reduces the number of free parameters in the gravitational sector from 2 to 1. If both $c_1(\mu)$ and $c_2(\mu)$ could be independently measured at any single scale $\mu$, one could test whether their relationship is consistent with the SCT prediction.

% ============================================================
\section{Scale-Dependent Ratio $c_1(\mu)/c_2(\mu)$}
% ============================================================

The ratio $c_1/c_2 = -1/3$ holds exactly only at $\mu = \Lambda$. At lower scales:
\begin{equation}
  \frac{c_1(\mu)}{c_2(\mu)} = \frac{f_4/(480\pi^2) - 2b_1\ln(\Lambda/\mu)}{-f_4/(160\pi^2) - 2b_2\ln(\Lambda/\mu)}
\end{equation}

For $\mu \ll \Lambda$, the logarithmic terms dominate and the ratio approaches:
\begin{equation}
  \frac{c_1(\mu)}{c_2(\mu)} \xrightarrow{\mu \ll \Lambda} \frac{b_1}{b_2}
\end{equation}
which is determined entirely by the matter content. For the SM:
\begin{equation}
  \frac{c_1(\mu)}{c_2(\mu)} \xrightarrow{\mu \ll \Lambda} \frac{b_1^{\text{SM}}}{b_2^{\text{SM}}} = \frac{363/10}{53/9} = \frac{3267}{530} \approx 6.16
\end{equation}

This dramatic departure from $-1/3$ at low energies means that the SCT prediction is most sharply testable at the highest accessible energies, i.e., closest to $\Lambda$.

% ============================================================
\section{Energy Scale and Regime}
% ============================================================

\begin{itemize}[leftmargin=2em]
  \item \textbf{Regime:} UV. The prediction specifies the running from $\mu = \Lambda$ (cutoff) downward.
  \item \textbf{Cutoff scale:} $\Lambda \sim M_P \approx 1.22 \times 10^{19}$~GeV (expected). The precise value depends on the spectral triple.
  \item \textbf{Where distinguishable from EFT:} At any scale where both $c_1$ and $c_2$ can be independently measured. The constraint is one relation between two quantities:
    \begin{equation}
      c_1(\mu) + \frac{1}{3}\,c_2(\mu) = -2\left(b_1 + \frac{1}{3}b_2\right)\ln\!\left(\frac{\Lambda}{\mu}\right)
      \label{eq:constraint}
    \end{equation}
    This linear combination is predicted as a function of $\mu/\Lambda$, with no additional free parameters (once $f_4$ is used to set the overall normalization).
  \item \textbf{RG validity:} The one-loop running~\eqref{eq:c1-running}--\eqref{eq:c2-running} is valid in the perturbative regime $c_i \ll 1/G\mu^2$. Higher-loop corrections become important near $\mu \sim M_P$.
\end{itemize}

% ============================================================
\section{Required Experimental Precision}
% ============================================================

\subsection{What Would Be Needed}

To test this prediction, one needs:
\begin{enumerate}
  \item Independent measurements of $c_1(\mu)$ and $c_2(\mu)$ at the same energy scale $\mu$.
  \item Sufficient precision to check the constraint~\eqref{eq:constraint}.
  \item Alternatively: measurement of $c_1(\mu)$ at two different scales to extract $b_1$ and $c_1(\Lambda)$ separately, then compare $c_1(\Lambda)$ with the SCT prediction.
\end{enumerate}

\subsection{Current Status}

\begin{center}
\renewcommand{\arraystretch}{1.2}
\begin{tabular}{lcc}
  \toprule
  \textbf{Observable} & \textbf{Sensitivity to $c_i$} & \textbf{Current precision} \\
  \midrule
  Graviton scattering (LHC) & $\propto c_i \cdot E^2/M_P^2$ & Not observed \\
  Black hole QNMs & $\propto c_i/(M_P^2 r_s^2)$ & $\delta\omega/\omega \sim 0.1$ \\
  Inflation (tensor tilt) & $\propto c_2 \cdot H^2/M_P^2$ & $r < 0.036$ (BICEP3) \\
  Neutron star mergers & $\propto c_i \cdot R_{\text{NS}}^{-4}$ & Not yet resolved \\
  \bottomrule
\end{tabular}
\end{center}

\noindent
In all cases, the sensitivity to $c_i$ is suppressed by at least $(E/M_P)^2$, making direct measurement impossible with current technology.

\subsection{Theoretical Constraints}

Certain combinations of $c_1, c_2$ are constrained by:
\begin{itemize}
  \item \textbf{Unitarity:} $c_2 < 0$ required to avoid the massive spin-2 ghost (Stelle, 1977). SCT satisfies this: $c_2^{\text{SCT}} = -f_4/(160\pi^2) < 0$ for $f_4 > 0$.
  \item \textbf{Positivity bounds:} Forward dispersion relations constrain certain linear combinations of $c_i$ (de Rham et al., 2017). Compatibility with SCT values is a nontrivial check.
  \item \textbf{Asymptotic safety:} If gravity is asymptotically safe, the UV fixed point determines $c_i(\Lambda_{\text{UV}})$. Comparing with SCT values could test or constrain both scenarios.
\end{itemize}

% ============================================================
\section{Remaining Uncertainty: The Moment $f_4$}
% ============================================================

The individual values of $c_1(\Lambda)$ and $c_2(\Lambda)$ depend on $f_4$. This quantity is determined by:
\begin{equation}
  f_4 = \int_0^\infty f(u)\,u\,du
\end{equation}
where $f$ is the cutoff function in the spectral action. The ratio $c_1/c_2 = -1/3$ is \emph{independent} of $f_4$, but the absolute scale of the running (and thus the constraint~\eqref{eq:constraint}) requires knowing $f_4$.

\noindent
\textbf{How $f_4$ is fixed in practice:}
\begin{itemize}
  \item By choosing a specific cutoff function $f$ (e.g., $f(u) = e^{-u}$ gives $f_4 = 1$).
  \item By matching the spectral action to a known physical quantity (e.g., the cosmological constant, the gravitational coupling, or a particle mass).
  \item In the noncommutative geometry approach (Connes--Chamseddine--Marcolli), $f_4$ is related to the Yukawa coupling matrix of the Standard Model.
\end{itemize}

% ============================================================
\section{Status}
% ============================================================

\begin{center}
\colorbox{untested!15}{\parbox{0.85\textwidth}{\centering
  \textbf{\color{untested}\Large UNTESTED}\\[6pt]
  No current experiment probes the running of curvature-squared couplings.\\[4pt]
  \textbf{Distinguishing power:} SCT reduces 2 free parameters to 1,\\
  imposing the constraint $c_1(\Lambda) + \frac{1}{3}c_2(\Lambda) = 0$.\\
  This is testable in principle if $c_1$ and $c_2$ are ever measured independently.\\[4pt]
  \textbf{Theoretical consistency:} The predicted sign $c_2 < 0$ is consistent with\\
  unitarity requirements for the massive spin-2 mode.\\
  The predicted initial conditions are compatible with known positivity bounds.
}}
\end{center}

% ============================================================
\section{Relation to Other SCT Predictions}
% ============================================================

\begin{itemize}[leftmargin=2em]
  \item \textbf{P-00 (Donoghue coefficient):} Model-independent; does not involve $c_1, c_2$.
  \item \textbf{P-01 (Fixed $c_1/c_2$ ratio):} Provides the initial condition at $\mu = \Lambda$ that is the starting point of the running described here.
  \item \textbf{P-03 (Nonlocal form factors):} Goes beyond the local running couplings to capture the full momentum-dependent structure of the spectral action.
\end{itemize}

\vspace{1em}
\noindent
\textit{This prediction card is part of the SCT Theory prediction catalog.}

\end{document}
